\subsection{Ext groups with simple modules}
\label{sec:standard-simple-extensions}


Using the projective resolutions constructed in
\Cref{sec:projective-resolutions-standard-modules}, we determine the Ext groups
$\Ext^q_{\Lambda_P}(\DD(T),\LL(C))$ for $q\geq0$ in terms of the
combinatorics of the intervals $T$ and $C$
(\Cref{thm:standard-simple-ext-formula}).
The corresponding simple--costandard statements follow by
poset duality.

\subsubsection{Standard--simple Ext groups}
\label{subsec:standard-simple-ext-groups}

Fix $T\in\Int(P)$ and $C\in\cU^\circ(T)$.
With $I_T(C)$ and
$\operatorname{comp}_J(C)$ as in
\eqref{eq:ITD-definition} and \eqref{eq:compJD-definition},
respectively, define
\begin{equation}
\label{eq:Gamma-definition}
\Gamma_T^-(C)
=
\bigl\{
J\subseteq I_T(C)
\bigm|
C\subsetneq\operatorname{comp}_J(C)
\bigr\}.
\end{equation}
If $J\in\Gamma_T^-(C)$ and $J'\subseteq J$, then
\eqref{eq:lifted-component-compatibility} shows that
$\operatorname{comp}_J(C)\subseteq \operatorname{comp}_{J'}(C)$. 
Hence $C\subsetneq\operatorname{comp}_{J'}(C)$ and $J'\in \Gamma_T^-(C)$. 
Thus
$\Gamma_T^-(C)$ forms a simplicial subcomplex of
$\Simp(I_T(C))$.

We first describe these Ext groups using the relative simplicial
cochain complex associated with this pair.

\begin{proposition}
\label{prop:Gamma-relative-cochain}
There is an isomorphism of cochain complexes
\begin{equation}
\label{eq:Gamma-relative-cochain-complex}
\Hom_{\Lambda_P}\bigl(\QQ_\bullet(T),\LL(C)\bigr)
\cong
C^\bullet
\bigl(\Simp(I_T(C)),\Gamma_T^-(C);\kk\bigr)[-1].
\end{equation}
Consequently, for every $n\geq1$,
\begin{equation}
\label{eq:Gamma-Ext-formula}
\Ext_{\Lambda_P}^n\bigl(\DD(T),\LL(C)\bigr)
\cong
\widetilde H^{n-2}\bigl(\Gamma_T^-(C);\kk\bigr).
\end{equation}
\end{proposition}

\begin{proof}
For every $A\in\Int(P)$,
\[
\Hom_{\Lambda_P}\bigl(\PP(A),\LL(C)\bigr)
\cong
e_A\LL(C)
\cong
\begin{cases}
\kk,&A=C,\\
0,&A\neq C.
\end{cases}
\]
Hence, after applying
$\Hom_{\Lambda_P}(-,\LL(C))$ to $\QQ_\bullet(T)$,
only the summands isomorphic to $\PP(C)$ remain.

We first determine which subsets $J$ index these summands.
By \eqref{eq:explicit-projective-terms}, a summand $\PP(C)$ occurs in
$\QQ_n(T)$ for a subset $J\subseteq[m]$ precisely when
\[
|J|=n
\qand
C\in\piZero(U_J).
\]
For the fixed interval $C$, the second condition is equivalent to
\[
J\subseteq I_T(C)
\qand
\operatorname{comp}_J(C)=C.
\]
Indeed, if $C\in\piZero(U_J)$, then $C\subseteq U_J$, so
\Cref{lem:unique-lift-admissible-support} gives
$J\subseteq I_T(C)$, and the unique component of $U_J$ containing
$C$ is $C$ itself. The converse follows immediately from
\eqref{eq:compJD-definition}.

Since
$C\subseteq\operatorname{comp}_J(C)$ for every
$J\subseteq I_T(C)$, \eqref{eq:Gamma-definition} now shows that the
surviving subsets $J$ with $|J|=n$ are precisely the
$(n-1)$-faces of $\Simp(I_T(C))$ that do not belong to
$\Gamma_T^-(C)$. Thus, for $n\geq1$, the summands on the two sides of
\eqref{eq:Gamma-relative-cochain-complex} are indexed by the same
subsets $J$. In degree zero, both sides vanish because
$\QQ_0(T)=\PP(T)$, $C\neq T$, and
$\varnothing\in\Gamma_T^-(C)$.
Taking $(-1)^n$ times the resulting identification in degree $n$
gives an isomorphism of graded vector spaces
\[
\Hom_{\Lambda_P}\bigl(\QQ_\bullet(T),\LL(C)\bigr)
\cong
C^\bullet
\bigl(\Simp(I_T(C)),\Gamma_T^-(C);\kk\bigr)[-1].
\]

It remains to check the differentials. Let
$K=\{j_0<\cdots<j_n\}\subseteq I_T(C)$ satisfy
$\operatorname{comp}_K(C)=C$, and put
$K_r=K\setminus\{j_r\}$.
The component associated with $K_r\subseteq K$ contributes after
applying $\Hom_{\Lambda_P}(-,\LL(C))$ precisely when
$\operatorname{comp}_{K_r}(C)=C$.
In that case,
\eqref{eq:explicit-projective-component-differential} gives the
corresponding component of $d_{n+1}$ as
\[
(-1)^r(j_{C,C})_*
=
(-1)^r\id_{\PP(C)}.
\]
Hence precomposition with $d_{n+1}$ gives the usual alternating
simplicial coboundary on the subsets $J$ that do not belong to
$\Gamma_T^-(C)$. The differential on the shifted cochain complex
$[-1]$ is the negative of the simplicial coboundary, and the factor
$(-1)^n$ in degree $n$ makes the degreewise identifications commute
with the differentials. This proves
\eqref{eq:Gamma-relative-cochain-complex}.

Since $\QQ_\bullet(T)$ is a projective resolution of $\DD(T)$ by
\Cref{thm:standard-projective-resolution}, taking cohomology in
\eqref{eq:Gamma-relative-cochain-complex} gives, for every $n\geq1$,
\[
\Ext_{\Lambda_P}^n\bigl(\DD(T),\LL(C)\bigr)
\cong
H^{n-1}
\bigl(\Simp(I_T(C)),\Gamma_T^-(C);\kk\bigr).
\]
Since $C\in\cU^\circ(T)$, the set $I_T(C)$ is nonempty, so
$\Simp(I_T(C))$ is a nonempty simplex and hence contractible.
The long exact sequence in relative cohomology
\cite[Section~3.1]{Hatcher02} therefore gives
\[
H^{n-1}
\bigl(\Simp(I_T(C)),\Gamma_T^-(C);\kk\bigr)
\cong
\widetilde H^{n-2}\bigl(\Gamma_T^-(C);\kk\bigr).
\]
This proves \eqref{eq:Gamma-Ext-formula}.
\end{proof}



Consider the \emph{lower boundary} of $C$ in $T$ defined by 
\begin{equation}
\label{eq:lower-boundary}
\partial_T^-C
:=
\bigl\{
x\in T\setminus C
\bigm|
x\lessdot_T c
\text{ for some }c\in C
\bigr\}. 
\end{equation}
We then define a simplicial complex 
\begin{equation}
\label{eq:boundary-Dowker-complex}
\cN_T^-(C)
:=
\bigl\{
F\subseteq\partial_T^-C
\bigm|
F\subseteq U_i
\text{ for some }i\in I_T(C)
\bigr\}.
\end{equation}


The following comparison is a consequence of Dowker duality \cite{Dowker52} for the
simplicial complexes $\Gamma_T^-(C)$ and $\cN_T^-(C)$.

\begin{proposition}
\label{prop:Gamma-Dowker}
There is a homotopy equivalence
\[
\Gamma_T^-(C)\simeq\cN_T^-(C).
\]
Consequently, for every $r\in\mathbb Z$,
\begin{equation}
\widetilde H^r\bigl(\Gamma_T^-(C);\kk\bigr)
\cong
\widetilde H^r\bigl(\cN_T^-(C);\kk\bigr).
\label{eq:Gamma-boundary-cohomology}
\end{equation}
\end{proposition}


\begin{proof}
Consider the finite relation
$\mathcal R\subseteq I_T(C)\times\partial_T^-C$ defined by
\[
i\mathrel{\mathcal R}x
:\Longleftrightarrow
x\in U_i.
\]

Let $J\subseteq I_T(C)$. We first observe that
\[
C\subsetneq\operatorname{comp}_J(C)
\Longleftrightarrow
U_J\cap\partial_T^-C\neq\varnothing.
\]
Indeed, if
$C\subsetneq\operatorname{comp}_J(C)$, connectedness gives
$x\in U_J\setminus C$ and $c\in C$ with $x<c$.
Choose such a pair with $[x,c]$ minimal.  Since $U_J$ is a relative
upset of $T$, every element strictly between $x$ and $c$ belongs to
$U_J$.  The minimality of $[x,c]$ therefore gives
$x\lessdot_T c$.  Hence
$x\in U_J\cap\partial_T^-C$.
Conversely, if $x\in U_J\cap\partial_T^-C$, then
$x\lessdot_T c$ for some $c\in C$.  Thus $x$ and $C$ belong to the same connected component of $U_J$, so
$C\subsetneq\operatorname{comp}_J(C)$.

Since $U_J=\bigcap_{i\in J}U_i$, it follows from
\eqref{eq:Gamma-definition} that
\[
J\in\Gamma_T^-(C)
\Longleftrightarrow
\text{there exists }x\in\partial_T^-C
\text{ such that }
i\mathrel{\mathcal R}x
\text{ for every }i\in J.
\]
On the other hand, by
\eqref{eq:boundary-Dowker-complex}, for
$F\subseteq\partial_T^-C$,
\[
F\in\cN_T^-(C)
\Longleftrightarrow
\text{there exists }i\in I_T(C)
\text{ such that }
i\mathrel{\mathcal R}x
\text{ for every }x\in F.
\]

Thus $\Gamma_T^-(C)$ and $\cN_T^-(C)$ are precisely the two
common-neighbor complexes associated with the finite relation
$\mathcal R$. The Bipartite Relation Theorem
\cite[Theorem~10.9]{Bjorner95} gives
\[
\Gamma_T^-(C)\simeq\cN_T^-(C).
\]
The isomorphisms in reduced cohomology follow.
\end{proof}


By \Cref{prop:Gamma-Dowker}, it remains to compute the reduced
cohomology of $\cN_T^-(C)$. This is determined by an extremality
condition on the lower boundary, which we record in the following
definition.


\begin{definition}
\label{def:lower-extremal-boundary-pairs}
For $T\in\Int(P)$, put $\partial_T^-T:=\varnothing$. 
For $C\in \cU(T)$, we say that
$C$ has \emph{extremal lower boundary} in $T$ if
\[
\partial_T^-C=\Min(T\setminus C).
\]
In this case, we call $(T,C)$ a \emph{lower extremal-boundary pair}
in $P$ and put
\[
\omega^-(T,C)
:=
\bigl|\partial_T^-C\bigr|
=
\bigl|\Min(T\setminus C)\bigr|.
\]
In particular, $T$ has extremal lower boundary in itself with $\omega^-(T,T)=0$.
We denote by $\Upsilon^-(P)$ the set of all lower extremal-boundary
pairs in $P$.
\end{definition}




\begin{proposition}
\label{prop:boundary-complex-topology}
For $T\in \Int(P)$ and $C\in \cU^\circ(T)$, there is a homotopy equivalence
\begin{equation}
\label{eq:boundary-complex-topology}
\cN_T^-(C)\simeq
\begin{cases}
\partial\Simp\bigl(\Min(T\setminus C)\bigr)
\simeq
\mathbb S^{\omega^-(T,C)-2},
&
(T,C)\in\Upsilon^-(P),
\\
\mathrm{pt},
&
(T,C)\notin\Upsilon^-(P).
\end{cases}
\end{equation}
Consequently, for every $r\in\mathbb Z$,
\begin{equation}
\label{eq:boundary-complex-cohomology}
\widetilde H^r\bigl(\cN_T^-(C);\kk\bigr)
\cong
\begin{cases}
\det\bigl(\Min(T\setminus C)\bigr),
&
(T,C)\in\Upsilon^-(P)
\qand
r=\omega^-(T,C)-2,
\\
0,
&
\text{otherwise}.
\end{cases}
\end{equation}
\end{proposition}

\begin{proof}
Let $T\in \Int(P)$ and $C\in \cU^\circ(T)$.
We first claim that
\begin{equation}
\label{eq:boundary-minimal-form}
\cN_T^-(C)
=
\bigl\{
F\subseteq\partial_T^-C
\bigm|
\Min(T\setminus C)\nsubseteq F
\bigr\}.
\end{equation}


Recall that
\[
i\in I_T(C)
\Longleftrightarrow
C\subseteq U_i.
\]
Hence, 
\begin{align}
\cN_T^-(C)
&=
\bigl\{
F\subseteq\partial_T^-C
\bigm|
F\subseteq U_i
\text{ for some }i\in I_T(C)
\bigr\}
\notag
\\
&=
\bigl\{
F\subseteq\partial_T^-C
\bigm|
C\cup F\subseteq U_i
\text{ for some }i\in[m]
\bigr\}
\notag
\\
&=
\bigl\{
F\subseteq\partial_T^-C
\bigm|
C\cup F\subseteq U
\text{ for some }U\in\cU^\circ(T)
\bigr\}.
\label{eq:boundary-proper-upset-characterization}
\end{align}

For $F\subseteq\partial_T^-C$, we now compare the condition in
\eqref{eq:boundary-proper-upset-characterization} with that in
\eqref{eq:boundary-minimal-form}. If
$\Min(T\setminus C)\subseteq F$,
then every relative upset containing $C\cup F$ is all of $T$.
Thus $F\notin\cN_T^-(C)$ by
\eqref{eq:boundary-proper-upset-characterization}.

Conversely, suppose that
$\Min(T\setminus C)\nsubseteq F$
and consider 
\[
U=(C\cup F)^{\uparrow_T}.
\]
Then $U$ is a proper relative upset of $T$ containing $C\cup F$. 
Moreover, $U$ is connected since $F\subseteq\partial_T^-C$ and $C$ is connected. Thus $U$ satisfies the condition in
\eqref{eq:boundary-proper-upset-characterization}, and therefore
$F\in\cN_T^-(C)$. This proves
\eqref{eq:boundary-minimal-form}.


Suppose first that $(T,C)\in\Upsilon^-(P)$. Then
$\partial_T^-C=\Min(T\setminus C)$, so
\eqref{eq:boundary-minimal-form} gives
\[
\cN_T^-(C)
=
\partial\Simp\bigl(\Min(T\setminus C)\bigr).
\]
Since
$\omega^-(T,C)=|\Min(T\setminus C)|$ by
\Cref{def:lower-extremal-boundary-pairs}, we obtain
\[
\cN_T^-(C)
\simeq
\mathbb S^{\omega^-(T,C)-2}.
\]


Suppose next that $(T,C)\notin\Upsilon^-(P)$. If
$\Min(T\setminus C)\nsubseteq\partial_T^-C$, then the condition in
\eqref{eq:boundary-minimal-form} is automatic. Hence $\cN_T^-(C)$ is
the full simplex on $\partial_T^-C$. Since $T$ is connected and
$C$ is a proper relative upset, $\partial_T^-C$ is nonempty, and
thus $\cN_T^-(C)$ is contractible.

Otherwise,
$\Min(T\setminus C)\subsetneq\partial_T^-C$, and
\eqref{eq:boundary-minimal-form} gives
\[
\cN_T^-(C)
=
\partial\Simp\bigl(\Min(T\setminus C)\bigr)
*
\Simp\bigl(
\partial_T^-C\setminus\Min(T\setminus C)
\bigr).
\]
The second factor is a simplex on a nonempty vertex set, so this join
is contractible. This proves
\eqref{eq:boundary-complex-topology}.

In the first case, the canonical identification of the top reduced
cohomology of the boundary of a simplex gives
\[
\widetilde H^{\omega^-(T,C)-2}
\bigl(\cN_T^-(C);\kk\bigr)
\cong
\det\bigl(\Min(T\setminus C)\bigr),
\]
and all other reduced cohomology groups vanish. In the remaining
cases, $\cN_T^-(C)$ is contractible. This proves
\eqref{eq:boundary-complex-cohomology}.
\end{proof}


The preceding results yield the main result of this subsection.

\begin{theorem}
\label{thm:standard-simple-ext-formula}
For $T,C\in\Int(P)$ and $p\geq0$,
\begin{equation}
\label{eq:standard-simple-Ext}
\Ext^p_{\Lambda_P}\bigl(\DD(T),\LL(C)\bigr)
\cong
\begin{cases}
\det\bigl(\Min(T\setminus C)\bigr),
&
(T,C)\in\Upsilon^-(P)
\qand
p=\omega^-(T,C),
\\
0,
&
\text{otherwise}.
\end{cases}
\end{equation}
\end{theorem}


\begin{proof}
For $p\geq1$, suppose first that $C\in\cU^\circ(T)$.
Then the formula follows directly from
\eqref{eq:Gamma-Ext-formula},
\eqref{eq:Gamma-boundary-cohomology}, and
\eqref{eq:boundary-complex-cohomology}.
If $C\notin\cU^\circ(T)$, the positive-degree terms of the projective
resolution $\QQ_\bullet(T)$ contain no summand $\PP(C)$, so the
corresponding Ext groups vanish.

For $p=0$, the simple top of $\DD(T)$ is $\LL(T)$. Hence
$\Hom_{\Lambda_P}(\DD(T),\LL(C))$ is one-dimensional if $C=T$ and
zero otherwise. Since $\omega^-(T,T)=0$, while $\omega^-(T,C)>0$ for
every lower extremal-boundary pair with $C\neq T$, this gives
\eqref{eq:standard-simple-Ext} in degree zero.
\end{proof}



\subsubsection{Simple--costandard Ext groups via poset duality}
\label{subsec:simple-costandard-ext-via-poset-duality}
Order reversal and $\kk$-linear duality also give an injective
coresolution of each costandard module $\NN(T)$ from the projective
resolution in \Cref{thm:standard-projective-resolution}, with maximal
proper relative downsets replacing maximal proper relative upsets.
Applying the same duality directly to the preceding standard--simple
calculation gives the corresponding simple--costandard Ext groups as
follows.


For $S,T\in \Int(P)$ with $S\in\cD(T)$, put
\[
\partial_T^+S
:=
\{x\in T\setminus S\mid s\lessdot_T x
\text{ for some }s\in S\}.
\]
We say that $S$ has \emph{extremal upper boundary} in $T$ if
\[
\partial_T^+S=\Max(T\setminus S),
\]
and in this case call $(S,T)$ an \emph{upper extremal-boundary pair}.
For such a pair, put
\[
\omega^+(S,T)
:=
|\partial_T^+S|
=
|\Max(T\setminus S)|.
\]
In particular, $T$ has extremal upper boundary in itself with
$\omega^+(T,T)=0$.
Let $\Upsilon^+(P)$ be the set of all upper extremal-boundary pairs.

Under poset duality, relative downsets in $T$ become relative upsets,
and the upper boundary computed in $P$ becomes the lower boundary computed
in $P^{\op}$. Thus
\[
(S,T)\in\Upsilon^+(P)
\Longleftrightarrow
(T,S)\in\Upsilon^-(P^{\op}).
\]

\begin{corollary}
\label{cor:simple-costandard-ext-formula}
For $S,T\in\Int(P)$ and $p\geq0$,
\begin{equation}
\label{eq:simple-costandard-Ext}
\Ext_{\Lambda_P}^p\bigl(\LL(S),\NN(T)\bigr)
\cong
\begin{cases}
\det\bigl(\Max(T\setminus S)\bigr),
&
(S,T)\in\Upsilon^+(P)
\qand
p=\omega^+(S,T),
\\
0,
&
\text{otherwise}.
\end{cases}
\end{equation}
\end{corollary}

\begin{proof}
The isomorphism
$\Lambda_{P^{\op}}\cong\Lambda_P^{\op}$ together with $\kk$-linear
duality gives a contravariant exact equivalence that exchanges
costandard modules for $P$ with standard modules for $P^{\op}$ and
preserves simple labels. We therefore have
\[
\Ext_{\Lambda_P}^p
\bigl(\LL_P(S),\NN_P(T)\bigr)
\cong
\Ext_{\Lambda_{P^{\op}}}^p
\bigl(\DD_{P^{\op}}(T),\LL_{P^{\op}}(S)\bigr).
\]
Applying \Cref{thm:standard-simple-ext-formula} to $P^{\op}$ and using
the preceding identification of upper extremal-boundary pairs with
lower extremal-boundary pairs gives
\eqref{eq:simple-costandard-Ext}.
\end{proof}

